\documentclass[10pt,a4paper]{amsart}
\usepackage[margin=19mm]{geometry}
\usepackage{amsmath,amssymb,mathtools}
\usepackage[scr=rsfs,cal=euler]{mathalfa}
\usepackage{xfrac}
\usepackage[unicode,hidelinks,bookmarks=false]{hyperref}
\newcommand{\ie}{i.e.\ }
\newcommand{\cf}{cf.\ }
\newcommand{\resp}{respectively\ }
\newcommand{\Subsection}{\subsection}
\providecommand{\nobreakdash}{\nobreak\hbox{-}}
\setlength{\emergencystretch}{2em}
\multlinegap0pt
\usepackage{amssymb}
\usepackage{enumerate}
\usepackage{float}
\newtheorem{lemma}{Lemma}
\title{Linear complementary dual \\quasi-cyclic codes of index 2}
\author{Kanat Abdukhalikov, Duy Ho, San Ling, Gyanendra K. Verma}
\date{Source: arXiv:2504.09126v3 (CC BY 4.0)}
\begin{document}
\maketitle

\noindent\emph{Source and licence.} Linear complementary dual \\quasi-cyclic codes of index 2. Version arXiv:2504.09126v3 (CC BY 4.0), distributed by arXiv under the Creative Commons Attribution 4.0 International licence, http://creativecommons.org/licenses/by/4.0/, as recorded in the arXiv licence metadata for this exact version and shown on the abstract page https://arxiv.org/abs/2504.09126v3. Full licence notice: \texttt{licenses/arxiv-2504.09126-CC-BY-4.0.txt}.\par\smallskip

\noindent\textit{Editorial scope.} Counted unit: the article's lemma s2 (source0.tex lines 602--603). The article's complete original proof of it is delivered with it (source0.tex lines 605--683, one \texttt{proof} environment of 79 lines). No mathematics is written, repaired or re-derived: the statement and the proof are the article's own lines, in the article's own notation, and no retained line is substituted or re-flowed.  Retained uncounted as the source-local context the proof needs: lines 472--473.  Nothing was omitted from any retained span, so the package carries no omission record.  No retained line contains a citation, so the wrapper carries no bibliography and no bibliography entry.  No retained line contains a figure environment, an image-inclusion command or drawing code, so no image is delivered and the package declares no asset dependency.  Editorial formatting only, in addition: an AMS \texttt{amsart} preamble that restates the article's own declarations and macros used by the retained lines, namely the environments \texttt{lemma} on separate counters, together with the standard packages those lines need.  The retained environments print in this document as Lemma 1 in order of appearance, whereas the article prints the same items with the numbering of its own full document.  The complete native source of the article as distributed in the arXiv e-print for this exact version is bundled unchanged as \texttt{sources/176419/source0.tex}.
\begin{lemma} \label{nIII} Assume that (I) and (II) hold. If (III) is not true, then $C$ is not Euclidean LCD. 
\end{lemma}

\begin{lemma} \label{s2}  If (I), (II), (III) and (IV) hold, then $C_j' \cap C_j''^{\perp_e} = \{\mathbf{0}\}$ and $C_j'' \cap C_j'^{\perp_e} = \{\mathbf{0}\}$  for all $1\le j \le p$.  
\end{lemma}

\begin{proof} Since  $g$ and $l$ are self-reciprocal, we have that
\begin{align*}
x^m-1   &= g \cdot l \cdot r_{11} \cdot t_{11}     \cdot  r_{22} \cdot  t_{22}.
\end{align*}
For each $j$, the irreducible polynomial $h_j$ divides exactly one of these six factors of $x^m-1$, which leads to the following six cases.


1. $h_j \mid g$. Since $g \mid g_{12}$, we also have that $h_j \mid g_{12}$.  Then  
$g_{11}(\xi^{v_j})=g_{12}(\xi^{v_j})=g_{22}(\xi^{v_j})=0$. Since $g$ is self-reciprocal, it follows that $h_j^* \mid g$. Similarly as before,
$\bar{g}_{11}(\xi^{v_j})=\bar{g}_{12}(\xi^{v_j})=\bar{g}_{22}(\xi^{v_j})=0$.
Hence $G_j'=G_j''=\{\mathbf{0}\}$, which implies that  $C_j' \cap C_j''^{\perp_e} = \{\mathbf{0}\}$ and $C_j'' \cap C_j'^{\perp_e} = \{\mathbf{0}\}$.


2. $h_j \mid l$. This condition implies that $h_j \nmid g$ and $h_j^* \nmid g$. Then $C_j'$ and $C_j''$ are 2-dimensional, which implies that $C_j'^{\perp_e}$ and $C_j''^{\perp_e}$ are $0$-dimensional. Then $C_j' \cap C_j''^{\perp_e} = \{\mathbf{0}\}$ and $C_j'' \cap C_j'^{\perp_e} = \{\mathbf{0}\}$.

3. $h_j \mid r_{11}$. This condition implies that $h_j \mid g_{11}$, $h_j \nmid g_{22}$, $h_j^* \mid g_{11}$ and $h_j^* \nmid g_{22}$. It follows that  
$$
G_j'=   \begin{bmatrix}
0 & g_{12}(\xi^{v_j}) \\
0 & g_{22}(\xi^{v_j}) 
\end{bmatrix} ,
G_j''=   \begin{bmatrix}
0 & \bar{g}_{12}(\xi^{v_j}) \\
0 & \bar{g}_{22}(\xi^{v_j}) 
\end{bmatrix}.
$$
It can then be readily checked that $C_j' \cap C_j''^{\perp_e} = \{\mathbf{0}\}$ and $C_j'' \cap C_j'^{\perp_e} = \{\mathbf{0}\}$.

4. $h_j \mid t_{11}$. We recall from the proof of Lemma \ref{nIII} that we can  rewrite 
\begin{align*}
x^m-1   &= \alpha \cdot g \cdot l \cdot r_{11} \cdot t_{11}^*    \cdot  r_{22} \cdot  t_{22}^*.
\end{align*}  
Then the condition $h_j \mid t_{11}$   implies that $h_j^* \mid t_{22}$. Hence $h_j \mid g_{11}$, $h_j \nmid g_{22}$, $h_j^* \nmid g_{11}$ and $h_j^* \mid g_{22}$.  
From condition (III), we have that $\gcd(t_{22},g_{12})=1$, and so $h_j^* \nmid g_{12}$.  
The matrices $G_j'$ and  $G_j''$ are of the following form
$$
G_j'=   \begin{bmatrix}
0 & g_{12}(\xi^{v_j}) \\
0 & g_{22}(\xi^{v_j}) 
\end{bmatrix} ,
G_j''=   \begin{bmatrix}
\bar{g}_{11}(\xi^{v_j}) & \bar{g}_{12}(\xi^{v_j}) \\
0 & 0
\end{bmatrix},
$$
where $g_{22}(\xi^{v_j})\ne0$, $\bar{g}_{11}(\xi^{v_j})\ne 0$, and $\bar{g}_{12}(\xi^{v_j})\ne0$.  Then $C_j' =\langle (0,1) \rangle$, $C_j'' =\langle (\bar{g}_{11}(\xi^{v_j}),\bar{g}_{12}(\xi^{v_j})) \rangle$, and the dual codes are 
$$
C_j'^{\perp_e}  =\langle (1,0) \rangle,C_j''^{\perp_e}  =\langle (-\bar{g}_{12}(\xi^{v_j}),\bar{g}_{11}(\xi^{v_j}))\rangle.
$$ 
It can then be readily checked that $C_j' \cap C_j''^{\perp_e} = \{\mathbf{0}\}$ and $C_j'' \cap C_j'^{\perp_e} = \{\mathbf{0}\}$.

5. $h_j \mid r_{22}$. This condition implies that $h_j \nmid g_{11}$, $h_j \mid g_{22}$, $h_j^* \nmid g_{11}$ and $h_j^* \mid g_{22}$. The matrices $G_j'$ and  $G_j''$ are of the following form
$$
G_j'=   \begin{bmatrix}
g_{11}(\xi^{v_j}) & g_{12}(\xi^{v_j}) \\
0 & 0
\end{bmatrix},
G_j''=   \begin{bmatrix}
\bar{g}_{11}(\xi^{v_j}) & \bar{g}_{12}(\xi^{v_j}) \\
0 & 0
\end{bmatrix},
$$
where $g_{11}(\xi^{v_j})\ne0$, and $\bar{g}_{11}(\xi^{v_j})\ne 0$. Then 
$$C_j' =\langle (g_{11}(\xi^{v_j}), g_{12}(\xi^{v_j})) \rangle , C_j'' =\langle (\bar{g}_{11}(\xi^{v_j}), \bar{g}_{12}(\xi^{v_j})) \rangle,$$
and the dual codes are 
$$
C_j'^{\perp_e}  =\langle (-g_{12}(\xi^{v_j}),g_{11}(\xi^{v_j}))\rangle, C_j''^{\perp_e}  =\langle (-\bar{g}_{12}(\xi^{v_j}),\bar{g}_{11}(\xi^{v_j}))\rangle.
$$  
From condition (IV), we have 
$$\gcd(r_{22}, g_{11}(x)\bar{g}_{11}(x)+g_{12}(x) \bar{g}_{12}(x))=1,$$ 
and since $h_j \mid r_{22}$, it follows that $h_j \nmid g_{11}(x)\bar{g}_{11}(x)+g_{12}(x) \bar{g}_{12}(x)$. Then
$$
g_{11}(\xi^{v_j})\bar{g}_{11}(\xi^{v_j})+g_{12}(\xi^{v_j}) \bar{g}_{12}(\xi^{v_j})) \ne 0.
$$
Hence $C_j' \cap C_j''^{\perp_e} = \{\mathbf{0}\}$ and $C_j'' \cap C_j'^{\perp_e} = \{\mathbf{0}\}$.

6. $h_j \mid t_{22}$. This case is similar to case 4.  \qedhere 

\end{proof}

\end{document}
